\documentclass{article}
\usepackage[UTF8]{ctex}
\usepackage{amsmath,amssymb,amsfonts}
\usepackage{enumitem}
\usepackage{geometry}
\geometry{a4paper, margin=1in}
% === 错题本排版优化（不参与题目标识） ===
\setlength{\parindent}{0pt}
\setlength{\parskip}{0.35em plus 0.10em minus 0.08em}
\setlength{\emergencystretch}{3em}
\setlist{itemsep=0.25em, topsep=0.40em, parsep=0pt}
\allowdisplaybreaks[2]
% === 错题本排版优化结束 ===
\begin{document}
\section*{南京师范大学2024年高等代数真题}
\section*{84．南京师范大学2024年高等代数试卷}

\subsection{题目1}
一．（15 分）设 \(f(x)\) 是整系数多项式，\(\frac{q}{p}\) 是 \(f(x)\) 的有理根，其中 \(p, q\) 是整数且 \((p, q)=1\) 。证明：对任意整数 \(m\) ，总有 \((p m-q) \mid f(m)\) 。

\subsection{题目7}
二．（20分）叙述并证明 Eisenstein 判别法则．

\subsection{题目8}
三．（15 分）计算 \(n\) 阶行列式

\[
D_{\mathbf{n}}=\left|\begin{array}{cccc}
a+x_{1} & a+x_{1}^{2} & \cdots & a+x_{1}^{n} \\
a+x_{2} & a+x_{2}^{2} & \cdots & a+x_{2}^{n} \\
\vdots & \vdots & & \vdots \\
a+x_{n} & a+x_{n}^{2} & \cdots & a+x_{n}^{n}
\end{array}\right| .
\]

\subsection{题目2}
四．（20 分）已知非齐次线性方程组

\[
\left\{\begin{array}{l}
x_{1}+x_{2}+x_{3}+x_{4}=-1 \\
4 x_{1}+3 x_{2}+5 x_{3}-2 x_{4}=-1 \\
a x_{1}+x_{2}+3 x_{3}+b x_{4}=1
\end{array}\right.
\]

有三个线性无关的解。

\begin{itemize}
  \item[（1）] 证明方程组的系数矩阵 \(A\) 的秩 \(r(A)=2\) ．
  \item[（2）] 求 \(a, b\) 的值及方程组的通解。
\end{itemize}

\subsection{题目3}
五．（20分）设 \(A\) 是 \(n\) 阶方阵，\(E+A\) 可逆，且 \(f(A)=(E-A)(E+A)^{-1}\) ．

\begin{itemize}
  \item[（1）] 证明：\((E+f(A))(E+A)=2 E\) 。
  \item[（2）] 证明：\(f(f(A))=A\) ．
\end{itemize}

\subsection{题目4}
六．（20 分）设 \(A\) 是 \(n\) 阶实对称矩阵，\(E\) 是 \(n\) 阶单位矩阵。证明：对充分小的正数 \(\varepsilon, E+\varepsilon A\) 为正定矩阵。

\subsection{题目5}
七．（20 分）设 \(W, W_{1}, W_{2}\) 是线性空间 \(V\) 的子空间，\(W_{1} \subseteq W, V=W_{1} \oplus W_{2}\) ．证明：

\[
\operatorname{dim} W=\operatorname{dim} W_{1}+\operatorname{dim}\left(W_{2} \cap W\right) .
\]

\subsection{题目6}
八．（20 分）设 \(V\) 是全体次数不超过 \(n\) 的实系数多项式再添上零多项式组成的实数域上的线性空间，定义 \(\sigma\) 是 \(V\) 上的线性变换：任给 \(f(x) \in V, \sigma(f(x))=x f^{\prime}(x)-f(x)\) ．

\begin{itemize}
  \item[（1）] 求 \(\sigma\) 的核 \(\sigma^{-1}(0)\) 及 \(\sigma\) 的值域 \(\sigma V\) 。
  \item[（2）] 证明：\(V=\sigma^{-1}(0) \oplus \sigma V\) ．
\end{itemize}
\end{document}
